\documentclass[../mathNotesPreamble]{subfiles}

\providecommand{\relscalefact}{1.4}
\begin{document}
\relscale{\relscalefact}
  \section{10.1: Trails, Paths, and Circuits}

  %TODO define walk, trail, path, closed walk, circuit, simple circuit (p679)
  \begin{defn*}
    \begin{itemize}
      \item
        Let $G$ be a graph, and let $v$ and $w$ be vertices in $G$.

      \item
        A \textbf{walk from $v$ to $w$} is a finite alternating sequence of adjacent vertices and edges of $G$. A walk has the form
          \[v_0\; e_1\; v_1\; e_2\; \dots\; v_{n-1}\; e_n\; v_n,\]
        where the $v$'s represent vertices, the $e$'s represent edges, $v_0=v$, $v_n=w$, and for each $i=1,2,\dots, n$, $v_{i-1}$ and $v_i$ are the endpoints of $e_i$. The \textbf{trivial walk from $v$ to $v$} consists of the single vertex $v$.

      \item
        A \textbf{trail from $v$ to $w$} is a walk from $v$ to $w$ that does not contain a repeated edge.

      \item
        A \textbf{path from $v$ to $w$} is a trail that does not contain a repeated vertex.

      \item
        A \textbf{closed walk} is a walk that starts and ends at the same vertex.

      \item
        A \textbf{circuit} is a closed walk that contains at least one edge and does not contain a repeated edge.

      \item
        A \textbf{simple circuit} is a circuit that does not have any other repeated vertex except the first and last.
    \end{itemize}
  \end{defn*}
  \pagebreak

  \begin{center}
    \renewcommand{\arraystretch}{1.5}
    \begin{tabular}{@{}m{30mm}*{2}{M{25mm}}*{2}{M{47.5mm}}@{}}
      \toprule
      & \textbf{Repeated Edge?}& \textbf{Repeated Vertex?}& \textbf{Starts and Ends at the Same Point?}& \textbf{Must Contain at Least One Edge?}\\\midrule
      \textbf{Walk}& allowed& allowed& allowed& no\\
      \textbf{Trail}& no& allowed& allowed& no\\
      \textbf{Path}& no& no& no& no\\
      \textbf{Closed walk}& allowed& allowed& yes& no\\
      \textbf{Circuit}& no& allowed& yes& yes\\
      \textbf{Simple} \newline\textbf{circuit}& no& first and last only& yes& yes\\
      \bottomrule
    \end{tabular}
  \end{center}
  \begin{ex*}
    Using the following graph:
    \vspace*{-1.5\baselineskip}

    \begin{center}
      \begin{tikzpicture}[
        node distance=20mm and 35mm,
        myNode/.style={circle, fill=black, inner sep=1.5pt},
        every loop/.style={looseness=40, -},
        label distance=2pt
      ]
        \node[myNode, label=45:$v_3$] (v3) at (0,0) {}
          edge [in=75,out=165,loop]
          node[above] {$e_4$} (v3);
        \node[myNode, right=of v3, label=0:$v_4$] (v4) {};
        \node[myNode, below=of v4, label=0:$v_5$] (v5) {};
        \node[myNode, below=of v3, label=-90:$v_6$] (v6) {};
        \node[draw=none] (v36) at ($(v3)!0.50!(v6)$) {};
        \node[myNode, left=of v36, label=-90:$v_2$] (v2) {};
        \node[myNode, left=of v2, label=-90:$v_1$] (v1) {};
        %%
        \draw[-] (v1) to [bend left=45] node[above] {$e_1$} (v2);
        \draw[-] (v1) to [bend right=45] node[below] {$e_2$} (v2);
        \draw[-] (v2) -- (v3) node[pos=0.5, above left] {$e_3$};
        \draw[-] (v3) -- (v4) node[pos=0.5, above] {$e_5$};
        \draw[-] (v3) -- (v5) node[pos=0.5, above right] {$e_6$};
        \draw[-] (v3) -- (v6) node[pos=0.5, left] {$e_7$};
        \draw[-] (v2) -- (v6) node[pos=0.5, below left] {$e_8$};
        \draw[-] (v6) -- (v5) node[pos=0.5, below] {$e_9$};
        \draw[-] (v5) -- (v4) node[pos=0.5, right] {$e_{10}$};
      \end{tikzpicture}
    \end{center}
    determine which of the following walks are trails, paths, circuits, or simple circuits.
    \begin{extasks}[after-item-skip=\stretch{1}](3)
      \task $v_1\,e_1\,v_2\,e_3\,e_4\,v_3\,e_5\,v_4$
      \task $e_1\,e_3\,e_5\,e_5\,e_6$
      \task $v_2\,v_3\,v_4\,v_5\,v_3\,v_6\,v_2$
      \task $v_2\,v_3\,v_4\,v_5\,v_6\,v_2$
      \task $v_1\,e_1\,v_2\,e_1\,v_1$
      \task $v_1$
    \end{extasks}
    \vspace*{\stretch{1}}
  \end{ex*}
  \pagebreak

  \begin{defn*}
    A graph $H$ is said to be a \textbf{subgraph} of a graph $G$ if, and only if, every vertex in $H$ is also a vertex in $G$, every edge in $H$ is also an edge in $G$, and every edge in $H$ has the same endpoints as it has in $G$.
  \end{defn*}
  \begin{ex*}
    Based on the following graph:
    \begin{center}
      \begin{tikzpicture}[
        node distance=20mm and 35mm,
        myNode/.style={circle, fill=black, inner sep=1.5pt},
        every loop/.style={looseness=40, -},
        label distance=2pt
      ]
        \node[myNode, label=90:$v_1$] (v1) at (0,0) {}
          edge [in=-45,out=45,loop]
          node[right] {$e_3$} (v1);
        \node[myNode, left=of v1, label=-90:$v_2$] (v2) {};
        %%
        \draw[-] (v1) to [bend right=45] node[above] {$e_1$} (v2);
        \draw[-] (v1) to [bend left=45] node[below] {$e_2$} (v2);
      \end{tikzpicture}
    \end{center}
    \begin{extasks}[after-item-skip=\stretch{1}](1)
      \task How many paths/trails/walks are there from $v_1$ to $v_2$?
%      \task How many trails are there from $v_1$ to $v_2$?
%      \task How many walks are there from $v_1$ to $v_2$?
      \task List all subgraphs.
    \end{extasks}
    \vspace*{\stretch{1}}
  \end{ex*}
  \pagebreak

  \begin{defn*}
    Let $G$ be a graph. Two \textbf{vertices $v$ and $w$ of $G$ are connected} if, and only if, there is a walk from $v$ to $w$. The \textbf{graph $G$ is connected} if, and only if, given \emph{any} two vertices $v$ and $w$ in $G$, there is a walk from $v$ to $w$.
  \end{defn*}

  \begin{thmBox*}[Lemma 10.1.1]
    Let $G$ be a graph.
    \begin{enumerate}
      \item If $G$ is connected, then any two distinct vertices of $G$ can be connected by a path.
      \item If vertices $v$ and $w$ are part of a circuit in $G$ and one edge is removed from the circuit, then there still exists a trail from $v$ to $w$ in $G$.
      \item If $G$ is connected and $G$ contains a circuit, then an edge of the circuit can be removed without disconnecting $G$.
    \end{enumerate}
  \end{thmBox*}
  \begin{defn*}
    A graph $H$ is a \textbf{connected component} of a graph $G$ if, and only if,
    \begin{enumerate}
      \item $H$ is a subgraph of $G$;
      \item $H$ is connected; and
      \item no connected subgraph of $G$ has $H$ as a subgraph and contains vertices or edges that are not in $H$.
    \end{enumerate}
  \end{defn*}
  \pagebreak

  \begin{ex*}
    Find all connected components of the following graph $G$
    \vspace*{\baselineskip}
    \begin{center}
      \begin{tikzpicture}[
        node distance=20mm and 35mm,
        myNode/.style={circle, fill=black, inner sep=1.5pt},
        every loop/.style={looseness=40, -},
        label distance=2pt
      ]
        \node[myNode, label=-90:$v_1$] (v1) at (0,0) {};
        \node[myNode, above=of v1, xshift= 15pt, label=90:$v_2$] (v2) {};
        \node[myNode, right=of v1, yshift=-20pt, label=0:$v_3$] (v3) {};
        \node[myNode, right=of v3, yshift= 30pt, label=-90:$v_4$] (v4) {};
        \node[myNode, right=of v4, yshift= 30pt, label=90:$v_5$] (v5) {};
        \node[myNode, right=of v4, yshift=-30pt, label=-90:$v_8$] (v8) {};
        \node[myNode, right=of v5, label=90:$v_6$] (v6) {};
        \node[myNode, right=of v8, label=-90:$v_7$] (v7) {};
        %%
        \draw[-] (v1) -- (v3) node[pos=0.5, below] {$e_1$};
        \draw[-] (v2) -- (v3) node[pos=0.5, right] {$e_2$};
        \draw[-] (v5) -- (v7) node[pos=0.25, above] {$e_3$};
        \draw[-] (v6) -- (v7) node[pos=0.5, right] {$e_4$};
        \draw[-] (v8) -- (v6) node[pos=0.25, below] {$e_5$};
      \end{tikzpicture}
    \end{center}
    \vspace*{\stretch{1}}
  \end{ex*}
  \pagebreak

  %TODO define Euler circuit (p684), theorem 10.1.2 (p684), theorem 10.1.3 (p685)
  \begin{defn*}
    Let $G$ be a graph. An \textbf{Euler circuit} for $G$ is a circuit that contains every vertex and every edge of $G$. That is, an Euler circuit for $G$ is a sequence of adjacent vertices and edges in $G$ that has at least one edge, starts and ends at the same vertex, uses every vertex of $G$ at least once, and uses every edge of $G$ exactly once.
  \end{defn*}

  \begin{thmBox*}[Theorem 10.1.2]
    If a graph has an Euler circuit, then every vertex of the graph has positive even degree.\\

    \textbf{Theorem 10.1.3}\newline
    If a graph $G$ is connected and the degree of every vertex of $G$ is a positive even integer, then $G$ has an Euler circuit.\\

    \textbf{Theorem 10.1.4}\newline
    A graph $G$ has an Euler circuit if, and only if, $G$ is connected and every vertex of $G$ has positive even degree.
  \end{thmBox*}

  \begin{defn*}
    Let $G$ be a graph, and let $v$ and $w$ be two distinct vertices of $G$. An \textbf{Euler trail from $v$ to $w$} is a sequence of adjacent edges and vertices that starts at $v$, ends at $w$, passes through every vertex of $G$ at least once, and traverses every edge of $G$ exactly once.
  \end{defn*}
  \pagebreak

  \begin{defn*}
    Given a graph $G$, a \textbf{Hamiltonian circuit} for $G$ is a simple circuit that includes every vertex of $G$. That is, a Hamiltonian circuit for $G$ is a sequence of adjacent vertices and distinct edges in which every vertex of $G$ appears exactly once, except for the first and the last, which are the same.
  \end{defn*}

  \begin{thmBox*}[Proposition 10.1.6]
    If a graph $G$ has a Hamiltonian circuit, then $G$ has a subgraph $H$ with the following properties:
    \begin{itemize}
      \item $H$ contains every vertex of $G$.
      \item $H$ is connected.
      \item $H$ has the same number of edges as vertices.
      \item Every vertex of $H$ has degree 2.
    \end{itemize}
  \end{thmBox*}
  %TODO traveling salesman problem (692)

  \begin{defn*}
    If $G$ is a simple graph, the \textbf{complement of $G$}, denoted $G'$, is obtained as follows: The vertex set of $G'$ is identical to the vertex set of $G$. However, two distinct vertices $v$ and $w$ of $G'$ are connected by an edge if, and only if, $v$ and $w$ are not connected by an edge in $G$.
  \end{defn*}

  \pagebreak
\end{document}
